\subsection{P007 --- A BIM-enabled digital twin framework for real-time indoor environment monitoring and visualization by integrating autonomous robotics, LiDAR-based 3D mobile mapping, IoT sensing, and indoor positioning technologies}
\label{paper:P007}
\textbf{Authors:} Xi Hu, Rayan H. Assaad\par
\textbf{Major Category:} Digital Twin\par
\textbf{Secondary Category:} Building Environment and IEQ, BIM and Semantic Information, Sensors, IoT and Data Integration\par
\textbf{Keywords:} Digital Twin, BIM, IoT, Indoor environmental monitoring, Indoor positioning, Robotics, LiDAR\par
\paragraph{主な主張}
BIM, IoT sensing, autonomous robot, LiDAR-based 3D mobile mapping, UWB indoor positioningを統合したreal-time indoor environment monitoring frameworkを二つの室内scenarioで検証した. 空気質の空間分布をBIM上で可視化し, UWB positioningはnearly centimeter-level precisionを示した. \cite{Hu:2024BIMDigitalTwin}
\paragraph{新規性}
従来別個に扱われていたBIM, DT, IoT, mobile robotic sensing, indoor positioningを単一frameworkへ統合し, 汚染源のlocalizationとvisualizationまで接続した点に特徴がある.
\paragraph{位置付け}
研究の中心が複数技術を統合したBIM-enabled digital twin frameworkであるためDigital Twinを主分類とする.
